\section{Classical inputs and conventions}\label{sec:background}

This section records the classical results used below, together with their
hypotheses and the conventions needed to apply them.
\subsection{Milnor lattices, spectra, and integral monodromy}

\begin{theorem}[Milnor]\label{bg:milnor}
Let $s\geq2$, and let $f:(\mathbb C^s,0)\to(\mathbb C,0)$ be a holomorphic germ with an isolated
critical point at $0$, with $f(0)=0$.  For sufficiently small $\varepsilon>0$, the link
$K_f=f^{-1}(0)\cap S^{2s-1}_\varepsilon$ is a smooth closed manifold of
dimension $2s-3$, and
\[
 f/|f|:S^{2s-1}_\varepsilon\setminus K_f\longrightarrow S^1
\]
is a locally trivial fibration.  Each fiber is the interior of its closure
$F_f$, a compact manifold with boundary $K_f$.  Both $F_f$ and its interior
are homotopy equivalent to a finite bouquet of $(s-1)$-spheres.
The link is $(s-3)$-connected.
\end{theorem}

See \cite[Theorems~4.8, 5.2, and~6.5, and Lemma~6.1]{Milnor1968}.
The free group $L_f=\widetilde H_{s-1}(F_f;\mathbb Z)$ is the
\emph{Milnor lattice}; for a one-variable power we use the reduced
zeroth homology of its finite Milnor fiber.  The unimodular integral Seifert form is
$S_f(x,y)=\operatorname{lk}(x,y^+)$, where the second cycle is pushed
positively to a neighboring page.
We use column vectors and fix the monodromy convention
\begin{equation}\label{eq:bg-monodromy}
 M_f=(-1)^sS_f^{-1}S_f^{\mathsf t}.
\end{equation}
This amounts to using the inverse of the monodromy in the convention
$S_f(Mx,y)=(-1)^sS_f(y,x)$; see
\cite[Section~6, Equations~(6.1)--(6.2)]{BalnojanHertling2019}.
Inverting monodromy throughout preserves every comparison made below.
An isomorphism of Seifert forms means integral congruence unless a different
coefficient field or ring is specified.

\begin{theorem}[Thom--Sebastiani--Sakamoto]\label{bg:ts}
Suppose $f$ and $g$ have isolated critical points in $a$ and $b$ variables,
respectively.  There is an isomorphism of integral lattices
\[
 L_{f+g}\cong L_f\otimes_{\mathbb Z}L_g
\]
under which
\[
 M_{f+g}\cong M_f\otimes M_g,
 \qquad
 S_{f+g}\cong(-1)^{ab}S_f\otimes S_g.
\]
The assertion includes one-variable powers, with reduced zeroth homology.
\end{theorem}

The tensor-product monodromy theorem originates with
Sebastiani and Thom~\cite{SebastianiThom1971}. The integral
Seifert-form formula is \cite[Theorem~2, p.~715]{Sakamoto1974};
the lattice and monodromy formulas are also recorded in
\cite[Section~6, Equations~(6.3)--(6.5)]{BalnojanHertling2019}.
Thus comparisons of sums with the same factor dimensions have the same
overall tensor sign.

A polynomial has \emph{normalized positive weights}
$w_1,\ldots,w_s\in\mathbb Q_{>0}$ if
$f(\lambda^{w_1}z_1,\ldots,\lambda^{w_s}z_s)=\lambda f(z)$, understood
after clearing denominators.  We use positive spectral exponents
$\beta\in(0,s)$, shifted by $1$ from the convention with spectrum in
$(-1,s-1)$.

\begin{theorem}[Weighted homogeneous spectrum]\label{bg:spectrum}
If $f$ is weighted homogeneous of normalized degree $1$ and has an isolated
critical point, its spectral polynomial is
\begin{equation}\label{eq:bg-spectrum}
 P_f(t)=\sum_{\beta\in\operatorname{Sp}^{+}(f)}t^\beta
       =\prod_{j=1}^s\frac{t^{w_j}-t}{1-t^{w_j}}.
\end{equation}
Exponents are counted with multiplicity.  Their classes modulo $\mathbb Z$
determine the eigenvalues of $M_f$ by exponentiation, with a fixed choice
of sign in the exponent.  The operator $M_f$ has finite order.
\end{theorem}

Here is the algebraic conversion from the usual spectrum formula to
\eqref{eq:bg-spectrum}.  Steenbrink's monomial formula
\cite[Section~1, Examples, p.~165]{Steenbrink1989} assigns exponent
$\sum_j(k_j+1)w_j$ to a monomial $z^k$ in a homogeneous basis of
$\mathbb C[z_1,\ldots,z_s]/(\partial f)$.
The derivatives have degrees $1-w_j$ and form a regular sequence:
their common zero set is the origin, so the ideal has height $s$ in the
Cohen--Macaulay polynomial ring
\cite[Tags~00ND and 02JN]{StacksProject}. The local regular-sequence
criterion applies at the homogeneous maximal ideal; positive grading
detects regularity there. Successive multiplication exact
sequences therefore give
\[
 \operatorname{Hilb}\bigl(\mathbb C[z]/(\partial f),t\bigr)
 =\prod_j\frac{1-t^{1-w_j}}{1-t^{w_j}}.
\]
Multiplication by $t^{\sum_jw_j}$ yields \eqref{eq:bg-spectrum}.
Finally, the weighted coordinate rotation
$z_j\mapsto e^{2\pi i w_j}z_j$ represents geometric monodromy; a common
multiple of the weight denominators kills this rotation.  Finite-order
operators in characteristic zero are semisimple.

\begin{theorem}[Real Seifert forms from the spectrum]\label{bg:real}
Let $f$ and $g$ be isolated weighted homogeneous singularities in the same
number of variables.  If their multisets of positive spectral exponents
agree modulo $2\mathbb Z$, their real Seifert forms are isomorphic.
\end{theorem}

For clarity, this is a consequence of the spectral-pair classification,
not an assertion about integral forms.  N\'emethi's theorem
\cite[Theorem~6.5]{Nemethi1995} and
\cite[Theorem~4.4(b)]{BalnojanHertling2019} say that spectral pairs modulo
$2\mathbb Z$ in the first coordinate, together with the dimension,
determine the real Seifert form.  In the present semisimple case the
nilpotent logarithm is zero.  Put $m=s-1$.  The actual cohomological weight
is $m$ away from eigenvalue $1$ and $m+1$ at eigenvalue $1$; the definition
of spectral pairs subtracts the latter shift, so every second coordinate
is $m$.  Thus this coordinate supplies no additional data.  Translation
$\beta\mapsto\beta-1$ respects equality modulo $2\mathbb Z$, and the
normalization and passage from cohomology to homology depend only on $s$.
More explicitly, the singularity construction gives a signed Steenbrink
polarized mixed Hodge structure, whose defining replacement of the
nilpotent logarithm by its negative leaves the data unchanged when that
logarithm is zero~\cite[Definition~6.1 and Remark~6.2]{BalnojanHertling2019}.
In that convention the normalized cohomological Seifert form is
$(-1)^{(m+1)(m+2)/2}$ times the dual homological Seifert form
\cite[Equation~(6.14)]{BalnojanHertling2019}.  The common sign and duality
preserve isomorphism.  This proves the stated consequence, including the
eigenvalue-$1$ part.

\begin{theorem}[Hertling--Mase]\label{bg:orlik}
Suppose $f$ is a sum of chains
\[
 x_1^{a_1}x_2+\cdots+x_{r-1}^{a_{r-1}}x_r+x_r^{a_r},
 \qquad a_i\geq2,
\]
including chains with $r=1$.  Write the characteristic polynomial of
$M_f$ as $\prod_d\Phi_d(t)^{c_d}$, where $\Phi_d$ is cyclotomic.
Then, as an integral module with operator $t=M_f$,
\begin{equation}\label{eq:bg-orlik}
 (L_f,M_f)\cong
 \bigoplus_{j=1}^{\max c_d}\mathbb Z[t]/\bigl(p_j(t)\bigr),
 \qquad p_j(t)=\prod_{c_d\geq j}\Phi_d(t).
\end{equation}
In particular, the eigenvalues with multiplicity determine this integral
monodromy module within the indicated class of singularities.
\end{theorem}

Orlik formulated the integral-monodromy decomposition
\cite{Orlik1972}; the chain calculation of Orlik and Randell
\cite{OrlikRandell1977} is a predecessor of the standard decomposition
proved by Hertling and Mase. The precise result used here is
\cite[Theorem~1.3(a), (c), (d), and Definition~2.6]{HertlingMase2022}.
Reversing the variables puts our chains in that paper's convention.
Each standard block has an integral cyclic generator, hence is exactly
the companion module displayed in \eqref{eq:bg-orlik}.  Inverting its
operator preserves cyclicity and its characteristic
polynomial, whose roots are invariant under inversion.  Thus \eqref{eq:bg-monodromy} causes no change
to the conclusion.  This theorem controls the lattice with monodromy;
it does not identify the integral Seifert form.

\subsection{Arithmetic inputs}

For a number field $F$, write $F_v$ for its completion at a place $v$.

\begin{theorem}[Hasse--Minkowski]\label{bg:hasse}
Two nondegenerate quadratic forms over a number field $F$ are isometric
over $F$ if and only if they are isometric over $F_v$ for every place
$v$, including the archimedean places.
\end{theorem}

We use the number-field form of the theorem; see
\cite[Theorem~25.64(c) and the following discussion]{Milne2022};
see also \cite[Chapter~VI]{OMeara1963}.

\begin{theorem}[Strong approximation]\label{bg:approximation}
Let $G$ be a connected, simply connected, absolutely almost simple
algebraic group over a number field $F$.  If
$\prod_{v\mid\infty}G(F_v)$ is noncompact, then $G(F)$ is dense in
$G(\mathbb A_{F,f})$, the group of finite adelic points.
\end{theorem}

See \cite[Theorem~A, p.~1140]{Platonov1969}, with the correction
in~\cite{Platonov1970}, and
\cite[Theorem~2.3, p.~281]{Rapinchuk2014}.
Here the finite adelic group is the restricted product of the $G(F_v)$
over finite places, with respect to integral compact open subgroups.
Consequently the density concerns approximation at finitely many
places together with integrality at all remaining places.
We apply the theorem to special unitary groups of nondegenerate
hermitian forms over a totally imaginary quadratic extension of a
totally real number field.  In dimension at least two these groups are
connected, simply connected, and absolutely almost simple; an
indefinite archimedean component makes their archimedean product
noncompact.  The construction below verifies these hypotheses for
every factor used.

\subsection{Embedded high-dimensional topology}

\begin{theorem}[Smale and Levine]\label{bg:classification}
A smooth closed manifold of dimension $d\geq5$ that is homotopy
equivalent to $S^d$ is homeomorphic to $S^d$.

Let $q\geq3$, and let $K_1,K_2\subset S^{2q+1}$ be smooth oriented
submanifolds homeomorphic to $S^{2q-1}$.  Suppose the knots are
\emph{simple}, meaning
\[
 \pi_1(S^{2q+1}\setminus K_j)\cong\mathbb Z,
 \qquad
 \pi_i(S^{2q+1}\setminus K_j)=0\quad(2\leq i\leq q-1).
\]
If their integral Seifert matrices are congruent, then they are
ambiently isotopic.  In particular their embedded sphere pairs are
homeomorphic.
\end{theorem}

The first assertion is the high-dimensional generalized Poincar\'e
theorem \cite[Theorem~A, p.~391]{Smale1961}.  The second follows from
\cite[Theorem~3, pp.~185--186]{Levine1970}: integral congruence is a
special case of the matrix equivalence in that theorem.  Levine's
knots are required to be homeomorphic to spheres, so exotic smooth
sphere structures are allowed.  Related formulations for Milnor
fibered links appear in \cite{Durfee1974},
\cite[pp.~698--699]{Saeki1999}, and
\cite[Section~2]{Neumann2017}.
The proof below checks the spherical and simplicity hypotheses
directly before applying this theorem.
